\section{The Generalized Perpetual Inventory Method: Recursive
	Construction and Calibration}
\label{app:A1:gpim_construction}

This appendix documents the recursive construction of U.S.
nonfinancial corporate capital stocks via the Generalized Perpetual
Inventory Method (GPIM) of \citet[App.~6.5]{Shaikh2016}. The
construction proceeds asset-by-asset for machinery and equipment
(ME) and nonresidential structures (NRC), recovering four distinct
quantities from BEA Fixed Assets data: the asset-specific implicit
price index $P_{j,t}$, the gross real stock $K^{G,R}_{j,t}$, the
net current-cost stock $K^{N,V}_{j,t}$, and the net-value schedule
$V_j(v)$. All series are expressed in constant 2017 dollars with
base-year normalization $P_{j,2017} = 100$.

\subsection{The Survival Function}
\label{app:A1:survival}

Physical survival of capital vintages follows a Weibull retirement
distribution. For asset class $j \in \{\text{ME}, \text{NRC}\}$,
the proportion of a vintage of age $v$ still in physical operation
is
\begin{equation}
	S_j(v) = \exp\!\left[
	-\left( \frac{v}{\lambda_j} \right)^{\alpha_j}
	\right],
	\quad v = 0, 1, \ldots, L_j,
	\label{eq:A1:survival}
\end{equation}
where $\alpha_j > 1$ is the Weibull shape parameter and the scale
parameter satisfies $\lambda_j = L_j / \Gamma(1 + 1/\alpha_j)$,
with $L_j$ denoting the mean service life. The shape parameter
$\alpha_j > 1$ implies a progressively increasing hazard rate:
capital goods remain fully operative in early years and are
withdrawn from service at an accelerating rate as they approach
the mean service life. This pattern is consistent with the
empirical evidence of \citet{Nomura2005}, who finds that nearly
half of 66 Japanese asset classes exhibit increasing hazard rates.

The calibrated parameters follow the BEA Fixed Assets methodology
\citep{Fraumeni1997}:
\begin{equation}
	L_{\text{ME}} = 14, \quad \alpha_{\text{ME}} = 1.7,
	\qquad
	L_{\text{NRC}} = 30, \quad \alpha_{\text{NRC}} = 1.6.
	\label{eq:A1:calibration_survival}
\end{equation}
The corresponding maximum ages are $L_{\text{ME}} = 14$ years and
$L_{\text{NRC}} = 30$ years, beyond which $S_j(v) = 0$ and the
vintage is fully retired.

\subsection{The Age-Price Schedule and Net-Value Function}
\label{app:A1:ageprice}

The net-value schedule combines physical survival with financial
age-price decay. For a vintage of age $v$, the net-value function
is
\begin{equation}
	V_j(v) = S_j(v) \cdot (1 - d_j)^v,
	\label{eq:A1:netvalue}
\end{equation}
where $d_j$ is the geometric declining-balance depreciation rate.
The GPIM separates the two dimensions that the BEA conflates in
its geometric depreciation convention: physical survival follows
the Weibull distribution \eqref{eq:A1:survival}, while financial
age-price decay follows the geometric profile $(1-d_j)^v$. This
separation is necessary to construct distinct capacity and
profitability measures \citep[App.~6.3]{Shaikh2016}.

The declining-balance rates are calibrated as
\begin{equation}
	d_{\text{ME}} = 0.110, \qquad d_{\text{NRC}} = 0.024,
	\label{eq:A1:calibration_decay}
\end{equation}
consistent with the geometric depreciation rates documented by
\citet{Fraumeni1997} for the corresponding asset classes.

\subsection{The SFC Implicit Price Recursion}
\label{app:A1:sfc_recursion}

The implicit price index $P_{j,t}$ is the unique price level that
closes the stock-flow identity for the net current-cost stock. In
each period $t$, the BEA-reported current-cost net stock must
satisfy
\begin{equation}
	K^{\text{NET,CC}}_{j,t}
	= I^{\text{NOM}}_{j,t}
	+ \frac{P_{j,t}}{100}
	\sum_{v=1}^{L_j} V_j(v)\, I^{\text{REAL}}_{j,t-v},
	\label{eq:A1:sfc_identity}
\end{equation}
where $I^{\text{NOM}}_{j,t}$ is current-year nominal gross
investment and $I^{\text{REAL}}_{j,t-v} = I^{\text{NOM}}_{j,t-v}
/ (P_{j,t-v}/100)$ is real investment of vintage $t-v$ deflated
by the previously recovered implicit price. The first term on the
right-hand side is the current-year investment vintage (age zero,
valued at the current price). The second term is the net-value-weighted
sum of surviving real vintages revalued at the current price level.

Solving \eqref{eq:A1:sfc_identity} for $P_{j,t}$ yields the
forward recursion
\begin{equation}
	P_{j,t} = 100 \cdot
	\frac{K^{\text{NET,CC}}_{j,t} - I^{\text{NOM}}_{j,t}}
	{\displaystyle\sum_{v=1}^{L_j} V_j(v)\,
		I^{\text{REAL}}_{j,t-v}}.
	\label{eq:A1:price_recursion}
\end{equation}
Given the initial price $P_{j,0}$ and the sequence of nominal
investment flows $\{I^{\text{NOM}}_{j,t}\}_{t=0}^{T}$, the
recursion \eqref{eq:A1:price_recursion} determines $P_{j,t}$
uniquely for each $t = 1, \ldots, T$. The price index is
``implicit'' because it is not directly observed: it is the unique
value that, applied to the survival-weighted real vintages,
reproduces the BEA's reported current-cost net stock.

\subsection{Initial Conditions and Base-Year Normalization}
\label{app:A1:initialization}

The recursion \eqref{eq:A1:price_recursion} requires an initial
price $P_{j,0}$ and initial real investment vintages. The base
year is 2017, at which $P_{j,2017} = 100$ by normalization. The
nominal investment series begins in 1901, sourced from BEA Fixed
Assets Table FAAt201 (nonresidential equipment and structures,
gross investment by asset class). The current-cost net stock
anchors $K^{\text{NET,CC}}_{j,t}$ are sourced from BEA Fixed
Assets Table FAAt401 (net stock by asset class, current-cost
basis).\footnote{Data retrieved via the BEA API
	(\texttt{DatasetName}: FixedAssets). Table FAAt201 provides
	nominal gross investment by asset; Table FAAt401 provides
	current-cost net stock by asset; Table FAAt402 provides
	current-cost depreciation by asset.}

The initialization proceeds backward from 2017. Given
$P_{j,2017} = 100$ and the observed nominal investment series,
the recursion \eqref{eq:A1:price_recursion} is iterated backward
to recover $P_{j,t}$ for $t < 2017$. Alternatively, the recursion
may be initialized at the earliest available year and iterated
forward; the two approaches yield identical results given the
BEA anchors.

\subsection{The Gross Real Stock}
\label{app:A1:gross_stock}

The capital stock used to estimate productive capacity is the gross
real stock, constructed as the survival-weighted sum of real
investment vintages:
\begin{equation}
	K^{G,R}_{j,t} = \sum_{v=0}^{L_j} S_j(v)\,
	I^{\text{REAL}}_{j,t-v}.
	\label{eq:A1:gross_stock}
\end{equation}
This measure relies exclusively on the physical survival schedule
$S_j(v)$ and excludes the age-price factor $(1-d_j)^v$. It
therefore measures the constant-price purchasing power embodied in
heterogeneous capital goods still in physical operation. Financial
write-downs, revaluation gains, and depreciation allowances do not
enter. The OECD manual confirms that the gross stock, valued at
``as new'' prices, ``is widely used as a broad indicator of the
productive capacity of a country'' \citep[31]{OECD2001}.

\subsection{The Warm-Up Period and Estimation Sample}
\label{app:A1:warmup}

The GPIM stock at any year $t$ requires investment vintages
extending back $L_j$ years. Because the nominal investment series
begins in 1901, the vintage summation for nonresidential structures
($L_{\text{NRC}} = 30$) is truncated for any year before 1931.
The stock values for 1901--1930 are therefore incomplete, omitting
pre-1901 vintages that would still be physically in service. The
estimation sample is restricted to 1931--2024, where the full
30-year vintage window is populated, yielding $T = 94$ annual
observations.

\subsection{Validation}
\label{app:A1:validation}

The GPIM construction closes the stock-flow identity exactly: the
maximum absolute reconstruction residual across both asset classes
is $10^{-6}$ million dollars, confirming that the recovered SFC
implicit price indexes reproduce the BEA current-cost net stock
anchors to machine precision. The BEA's Fisher quantity index is
retained solely as a validation comparison and does not enter the
GPIM baseline. The sectoral scope is restricted to nonfinancial
corporate (NFC) business, excluding non-profit institutions serving
households (NPISH), government, and owner-occupied housing.

\subsection{Data Sources}
\label{app:A1:data}

Table~\ref{tab:A1:data_sources} summarizes the BEA Fixed Assets
tables used in the GPIM construction.

\begin{table}[htbp]
	\centering
	\caption{BEA Fixed Assets tables used in GPIM construction}
	\label{tab:A1:data_sources}
	\small
	\begin{tabular}{@{}lll@{}}
		\toprule
		\textbf{Table} & \textbf{Content} & \textbf{Use} \\
		\midrule
		FAAt201 & Nominal gross investment by asset & $I^{\text{NOM}}_{j,t}$ \\
		FAAt401 & Current-cost net stock by asset & $K^{\text{NET,CC}}_{j,t}$ \\
		FAAt402 & Current-cost depreciation by asset & Validation \\
		\bottomrule
	\end{tabular}
	\begin{minipage}{\textwidth}
		\footnotesize
		\textit{Notes:} All tables are from the BEA Fixed Assets Accounts,
		retrieved via the BEA API (\texttt{DatasetName}: FixedAssets).
		Asset classes: ME = nonresidential equipment; NRC = nonresidential
		structures. Sectoral scope: nonfinancial corporate business.
	\end{minipage}
\end{table}